\documentclass[12pt]{article}
\usepackage[utf8]{inputenc}
\usepackage[LGR, T1]{fontenc}
\usepackage[greek.ancient,english]{babel}
\usepackage{alphabeta}

\usepackage[letterpaper, margin = 1in]{geometry}
\usepackage{fancyhdr}
\usepackage{titling}
\usepackage{titlesec}
\usepackage{hyperref}
\usepackage{lastpage}
\usepackage{amssymb}
\usepackage{setspace}
\usepackage[table]{xcolor}
\usepackage{tabularx}
\usepackage{array}
\usepackage{multirow}
\usepackage{arydshln}
\definecolor{graycol}{gray}{0.92}
\definecolor{darkergray}{gray}{0.75}

\makeatletter
\newcommand{\hangpara}[2]{\hangindent#1\hangafter#2\noindent}
\newenvironment{hangparas}[2]{\setlength{\parindent}{\z@}%
  \everypar={\hangpara{#1}{#2}}}{\par}
\makeatother

\providecommand{\medium}{\normalsize}

\title{Lysistrata's Unproductive Wool-Working and Its Origin in Penelope}
\author{Ian Thomas White}
\date{}

\makeatletter
\renewcommand{\maketitle}
{
\begin{flushleft}
\textbf{{\large\@title}}
\newline
\@author
\end{flushleft}
}
\makeatother

\titleformat*{\section}{\large\bfseries}
\titleformat*{\subsection}{\medium\bfseries}
\titleformat*{\subsubsection}{\small\bfseries}

\pagestyle{fancy}
\fancyhf{}
\setlength{\headheight}{14.5pt}
\lhead{Ian Thomas White}
\chead{Lysistrata's Unproductive Wool-Working}
\rhead{\thepage \ of \pageref{LastPage}}

\doublespacing
\widowpenalty=100000
\clubpenalty=100000

\begin{document}

\thispagestyle{empty}

\begin{singlespacing}
\maketitle
\tableofcontents
\end{singlespacing}

\newpage

\section{Introduction}

The wool-working metaphor in Aristophanes' \textit{Lysistrata}\footnote{All references to the Greek are from Henderson's edition: Aristophanes, \textit{Lysistrata}, ed. Jeffrey Henderson (Oxford: Clarendon Press, 1987). Translations are my own unless otherwise noted, in-text citations will be followed by line numbers.} is often remarked upon but remains relatively under-analyzed. Despite its centrality and marked elaborateness, the speech is often understood only insofar as it exemplifies some intrusion of a ``womanly'' domain into some ``manly'' domain.\footnote{E.g., Jeffrey Henderson, ``Commentary'', in \textit{Lysistrata}, ed. Jeffrey Henderson (Oxford: Clarendon Press, 1987), 141; Christopher A. Faraone, ``Salvation and Female Heroics in the Parodos of Aristophanes' Lysistrata'', \textit{The Journal of Hellenic Studies} 117, (1997), 58; Helene P. Foley, ``The `Female Intruder' Reconsidered: Women in Aristophanes' Lysistrata and Ecclesiazusae'', \textit{Classical Philology} 77, no. 1 (Jan. 1982), 7-9; Richard P. Martin, `` Fire on the Mountain: ``Lysistrata'' and the Lemnian Women'', \textit{Classical Antiquity} 6, no. 1 (Apr. 1987), 86; Laura McClure, ``Courtesans Reconsidered: Women in Aristophanes' Lysistrata'', \textit{Eugesta} 5 (2015), 72; John Vaio, ``The Manipulation of Theme and Action in Aristophanes' \textit{Lysistrata}'', \textit{Greek, Roman, and Byzantine Studies} 14 (Winter 1973), 373-376.} Sometimes scholars glean from the speech superficial political implications.\footnote{Tom Hawkins, ``Seducing a Misanthrope: timon the Philogynist in Aristophanes' Lysistrata'', \textit{Greek, Roman, and Byzantine Studies} 42 (2001), 157-158; Wilfred E. Major, \textit{The Court of Comedy: Aristophanes, Rhetoric and Democracy in Fifth-Century Athens} (Columbus: Ohio State University Press, 2013), 138; H.D. Westlake, ``The `Lysistrata' and the War'', \textit{Phoenix} 34, no. 1 (Spring 1980), 43} The crafty weaver Penelope is also sometimes blithely invoked, and usually in order to comment on Lysistrata's qualities as a ``powerful'' woman.\footnote{Matthew Dillon, ``The Lysistrata as a Post-Deceleian Peace Play'', \textit{Transactions of the American Philological Association} 117 (1987), 101; James McGlew, \textit{Citizens on Stage: Comedy and Political Culture in the Athenian Democracy} (Ann Arbor: University of Michigan Press, 2002), 161; Sarah Culpepper Stroup, ``Designing Women: Aristophanes' Lysistrata and the `Hetairization' of the Greek Wife'', \textit{Arethusa} 37, no. 1 (Winter 2004), 51.} The potential links between Lysistrata and Penelope in the domain of wool and weaving are largely left untouched, despite the pair's apparent superficial resemblances therein.\footnote{Ionna Papadopoulou-Belmehdi, ``Greek Weaving or the Feminine in Antithesis'', \textit{Diogenes} 42/3, no. 167 (Fall 1994), 44 and Lisa Pace Vetter, \textit{``Women's Work'' as Political Art} (Lanham, MD: Lexington Books, 2005), 63 link them insofar as both are part of a larger weaving theme in Greek literature.} The present paper attempts to show the textual dependence of \textit{Lysistrata's} wool-working scene\footnote{The ``scene'' being the wool-working metaphor speech (574-586) and its larger context in which the Proboulos questions Lysistrata (565-613).} on the descriptions in the \textit{Odyssey}\footnote{All references to the Greek and translation unless otherwise noted are from the Loeb edition: Homer, \textit{Odyssey}, vols. I-II, trans. A.T. Murray, LCL 104-105 (Cambridge, MA: Harvard University Press, 1919). Penelope's speech occurs three times (ii.96-102, xix.141-147, xxiv.131-137) and is exactly the same in each case. The immediately surrounding text in each case differs, but the scenes are essentially the same. I will provide only the first of a set of identical citations, and in-text citations will be followed by the book and line number.} of Penelope's weaving trick, a dependence which manifests in three ways: 1) simple textual and contextual links; 2) a (tenuous) metrical link; 3) a structural link, characterized by a potentially unending back-and-forth structural similarity. The political implications of this last structure and the Penelope-Lysistrata connection will also be explored.

\section{Textual and Contextual Links}

Athena, rather than Penelope, is often invoked as the model of choice for Lysistrata's character. Even Henderson goes so far as to say that ``Lysistrata finds her closest analogue in Athena herself''.\footnote{Jeffrey Henderson, ``Introduction'', in \textit{Lysistrata}, ed. Jeffrey Henderson (Oxford: Clarendon Press, 1987), xxxviii; Faraone, ``Salvation'', 58-59; Foley, ``Intruder'', 9-10} This comparison hinges on Athena's and Lysistrata's virginity and their mixtures of domestic and martial qualities.\footnote{Cf. Papadopoulou-Belmehdi, ``Greek Weaving'', 41-46.} In terms of weaving, the comparison more precisely hinges on Athena being the goddess of weaving and that fact's attendant cult practices, most notably the presentation of the \textit{peplos} (πέπλος) at the Panathenaia.\footnote{Henderson, ``Commentary'', 141. Whether or not the Athena connection is routed through the priestess Lysimache is outside the scope of this paper. Cf. Henderson, \textit{Lysistrata}, xxxviii-xli; Faraone, ``Salvation'', 58-59} Crucially, however, the garment that Lysistrata describes in her speech isn't a \textit{peplos} but a \textit{chlaina} (χλαῖνα, 586), the former being ``too delicate to serve the purposes she [Lysistrata] has here (dressing the male Demos in a stout garment)''.\footnote{Foley, ``Intruder'', 9 n. 20.} A direct reference to the \textit{peplos}\-based cult occurs when Women's Chorus claims to have been \textit{arrhephoroi} (641). The \textit{arrhephoroi}, however, ``supervise[d] the even younger girls who wove the peplos'',\footnote{Henderson, ``Commentary'', 155} and no mention of the actual weaving of the \textit{peplos} is mentioned here.

I don't mean to strictly deny any connection between Athena and Lysistrata; rather, the connection seems to be one of general persona rather than there being a direct link that's based on the weaving motif. No substantial weaving-based link, as far as I know, in the secondary scholarship has been made between Lysistrata and any of the other many weavers in Greek literature.\footnote{Even in Papadopoulou-Belmehdi, ``Greek Weaving'', which features several exemplars of the motif, Lysistrata is (among the pre-\textit{Lysistrata} texts) linked only to Penelope.} In addition to Penelope's already paradigm status as a literary weaver,\footnote{``[T]he most famous instance of this theme [endless weaving] is, of course, Penelope, simultaneously the most famous and the most obscure weaver in Greek mythology'', Papadopoulou-Belmehdi, ``Greek Weaving'', 42.} several positive and more specific links between Lysistrata's and Penelope's situations obtain.

First, the basic setup of the scene in which Lysistrata's wool-working speech occurs resembles important aspects of Penelope's spatial situation. Both women have as their ``base of operations'' secluded, exclusive, and gendered locations: Penelope's ὑπερῷον (e.g., i.328) and Lysistrata's freshly conquered acropolis.\footnote{Of course, Penelope's location isn't marked, and Lysistrata's command of the acropolis is dependent on the plot. Even if ``it is thus no accident that the women choose the Acropolis as their base of operations'', Faraone, ``Salvation'', 58, the transformation of the Acropolis into a gendered and cordoned zone is marked; i.e., the play makes a non-\textit{hyperoon}\-like setting into a \textit{hyperoon}\-like setting. Cf. McClure, ``Courtesans'', 62. This doesn't prove the \textit{hyperoon} refers to that of Penelope, but the settings are at least made to be similar. The seizure has, like other elements of the play. been traced to Athena, though, too (also e.g., Faraone, ``Salvation'', 58).} In these settings, Penelope and Lysistrata are singular leaders surrounded by subordinates.\footnote{There may be a link between Lysistrata's use of κάθημαι and Penelope as a paradigm case of the same word, McClure, ``Courtesans'', 62.} The subordinates, moreover, are depicted as potentially rebellious; Penelope's plan is discovered by the suitors when one of her maids snitched (τις ἔειπε γυναικῶν, ἣ σάφα ᾔδη, ii.108), and Lysistrata has to upbraid several women from deserting (728-780). On the outside of the female-controlled space is a male-gendered set of attackers. While Penelope's suitors are young in comparison to the play's chorus of old men, both groups of attackers are large sets of unruly men with a self-important leader, Antinous and the Proboulos, respectively.\footnote{The Proboulos exclaims his status (ὢν ἐγὼ πρόβουλος, 421) and needs money for ships (ἐκπορίσας ὅπως / κωπῆς ἔσονται, τἀργυρίου νυνὶ δέον, 421-422). Similarly, Antinous is one of two leaders of the suitors (ἀρχοὶ μνηστήρων, iv.629) and is described, e.g., as having a ``proud heart'' (θυμὸς ἀγήνωρ, iv.658); he also commands his fellows to give him a ship (δότε νῆα θοὴν, iv.669).} For extended communication to occur between each set of two camps, women must ``come down'' (καταβαίνω, e.g., i.328 in the \textit{Odyssey} and 864 in \textit{Lysistrata}) to men.

Second, both texts share the use of one particular word (in various forms) to describe both war and wool-working. The war in the \textit{Odyssey} is described as ``clewing up'' (πόλεμον τολυπεύειν, i.238, iv.490, xiv.368, xxiv.95), and Penelope self-describes her weaving actions as ``clewing up tricks'' (δόλους τολυπεύω, xix.137).\footnote{The actions of the women in Lysistrata are also briefly described by the Men's chorus as ``leading on with bait'' (ἐξεπαίρουσιν δόλῳ, 622).} Similarly, one part of Lysistrata's solution for domestic problems is to ``make a big clew [of wool]'' (ποιήσαι τολύπην μεγάλην, 585-586), a periphrasis quickly parroted by the Proboulos' ``clewing up'' (τολυπεύειν, 587). The Proboulos explicitly mentions ``clewing up'' (τολυπεύειν, 587) as an activity separate from the ``war'' (πόλεμος, 588), and Lysistrata doesn't try to salvage the connection while defending herself. While clew-making pairs war and wool-working in both texts, the felicity of that pairing is explicitly made dubious in \textit{Lysistrata}.

Finally, while Lysistrata's hypothetically weaved \textit{chlaina} (586) is a garment functionally distinct from Athena's \textit{peplos}, it is not so distinct from Penelope's \textit{pharos} (φᾶρος, ii.97). Both garments would have been worn over the \textit{chiton,} and the \textit{chlaina} ``is also called the \textit{pharos} by Homer''.\footnote{\textit{LSJ, s.v.} ``χλαῖνα'', last accessed November 2018, \url{http://www.perseus.tufts.edu/hopper/}; \textit{LSJ, s.v.} ``φᾶρος'', last accessed November 2018, \url{http://www.perseus.tufts.edu/hopper/}.} Of course, the \textit{pharos} made by Penelope is a funeral shroud (ταφήιος, ii.97), which wouldn't directly fit the sense in Lysistrata. However, the speech in \textit{Lysistrata} is paired with a funereal context in which Lysistrata mock prepares the Proboulos for his own burial (599-607). A quick sequence occurs from the speech, in which the Proboulos is ``taught to operate woolworking equipment'',\footnote{And he is ``taught to operate woolworking equipment'', Henderson, ``Commentary'', 129.} to a discussion of men and women growing old, and finally to the mock burial. Old age, moreover, is explicitly tied to women wasting away in their weaving chambers. Previously, though, before the wool-working speech, Lysistrata and a group of old women had already outfitted the Proboulos in a dress and with sewing materials (531-538). Taking the age-weaving discussion and the Proboulos as weaver together, a strong link is made between the hypothetical \textit{chlaina} and the Proboulos' death---the Proboulos has the materials and the knowhow to make his own shroud.\footnote{McClure, ``Courtesans'', 72 exploits the potential nuptial connotations of the \textit{chlaina}, positioning Lysistrata as the ``bride'' of the male demos. McClure's analysis, though, neglects the fact that the Proboulos is given the sewing materials and that the speech is constructed exclusively around an irreal impersonal verb with a set of accusative-subject-less infinitives. The most likely agent of these infinitives is some ambiguous ``we'', which could comprise any subset of Lysistrata, the other women, the Proboulos, and the entire polis.}

The purpose for Penelope's \textit{pharos}, then, seems indirectly referenced via the unfulfilled purpose of Lysistrata's \textit{chlaina}, namely, covering the Proboulos upon his death. This link better explains the suddenness of Lysistrata's harsh turn compared to her merely being fed up with the Proboulos' vulgarity. Of course, the Penelope-Lysistrata analogy switches from Antinous-Proboulos to Laertes-Proboulos, but this switch is paired with a mismatch between \textit{chlaina} designer/expert (Lysistrata) and \textit{chlaina} maker (the Proboulos). The supposedly dead Proboulos, also, precisely doesn't receive/produce the hypothetical shroud, unlike Laertes/Penelope. The scene, then, portrays a mismatch between design and product that fails to bring about the implicit goal. This mismatch, while seemingly funded by the logic of Penelope's situation, hints at the fallow nature of Lysistrata's hoped-for \textit{chlaina}\-covered demos, as well the general ineffectiveness of the women's (hypothetical) weaving designs.\footnote{It might be argued that the scene is more Proboulos-directed, i.e., that the women put the materials under the Proboulos' care to show his ineffectiveness at producing a weaving product (and thus, by extension, his inability to procure peace). Still, though, the Proboulos' burial is curiously incomplete, without a shroud, especially given the \textit{chlaina}\-based boasting of Lysistrata just previous. Furthermore, to foreshadow my own later arguments, the production, the actual weaving process, of the \textit{chlaina} itself is absent compared to the detailed description of wool-working preparation.}

These circumstantial links between Penelope and Lysistrata, if not conclusive, provide a strong grounding for the two characters' connection.

\section{A Tenuous Metrical Link}

The meter of Penelope's direct speech request to the suitors is epic dactylic hexameter, and the meter of Lysistrata's wool-working speech is anapestic tetrameter catalectic.\footnote{Henderson, ``Commentary'', 68.} The scansion of the latter might subtly refer to the scansion of the former in one particular way. Namely, the syllable structure of every fifth foot and of every second-to-last foot of each line of both texts' sections (the fifth foot in the \textit{Odyssey} in both cases and the fifth and seventh feet in \textit{Lysistrata}\footnote{And the fifth foot in \textit{Lysistrata} could be considered the first sub-foot of the third, or second-to-last foot, in the tetrameter structure. I.e., all the relevant feet are ``second-to-last'' feet}) are all perfectly, internally consistent with each other. Moreover, the second-to-last feet\footnote{The proper second-to-last feet.} in both texts are all the base form syllable structure of their respective meters. Thus, every fifth foot in Penelope's speech is a base-form dactyl, and the fifth foot is the only foot in this section which doesn't vary in syllable structure. Every fifth foot in Lysistrata's speech is a spondee, and every seventh foot is a base-form anapest, and all the other feet in the section vary.

In this larger anapestic tetrameter section of \textit{Lysistrata}, lines with a fifth-foot spondee and seventh-foot anapest admittedly seem fairly common, which makes the argument difficult to accept. To partially save it, I'll note that the fifth foot of the line directly before the speech, 573, is an anapest. Also, the fifth foot of the line that abruptly ends this section (when Lysistrata initiates the Proboulos' funeral), 593, is also an anapest. Thus, at least in terms of the fifth and the larger second-to-last foot of the tetrameter line, the speech is weakly framed. No such frame immediately occurs for the Penelope speech, in terms of the fifth foot, but this is a less significant obstacle. In any case, I provide the scansions below.\footnote{I use ``---'' for long syllable and ``x'' for short syllable. Vertical lines indicate caesura, primary or secondary. The ``()'' in 574 indicates that the first two feet are spoken by a preceding character (the Proboulos). In the seventh foot of line 585, the penultimate in ποιῆσαι is short, as is common in Attic poetry. \textit{LSJ}, \textit{s.v.} ``ποιέω'',  last accessed November 2018, \url{http://www.perseus.tufts.edu/hopper/}.}


\begin{table}[htbp]
\centering
\begin{singlespacing}
\textbf{Scansion of Lysistrata's wool-working speech (574--586) + frame (573, 594)}

\vspace{6pt}
\begin{tabular}{r | c c c c >{\columncolor{graycol}}c c >{\columncolor{graycol}}c c |}
573 & --- --- & x x --- & \; --- --- & \; x x --- & \cellcolor{darkergray} | x x --- & \; --- --- & x x --- & \multicolumn{1}{c}{x} \\
\hline
574 & (--- --- & x x ---) & | --- --- & \; x x --- & | --- --- & \; x x --- & x x --- & --- \\
575 & --- --- & --- --- & \; --- --- & \; --- --- & | --- --- & \; x x --- & x x --- & --- \\
576 & --- --- & --- --- & \; --- --- & \; --- --- & | --- --- & \; x x --- & x x --- & --- \\
577 & --- --- & x x --- & \; x x --- & \; --- --- & | --- --- & \; --- --- & x x --- & --- \\
578 & x x --- & --- --- & \; x x --- & \; --- --- & | --- --- & \; x x --- & x x --- & --- \\
579 & --- --- & --- --- & \; --- x x & \; --- --- & | --- --- & \; --- --- & x x --- & --- \\
580 & x x --- & --- --- & | --- x x & \; --- --- & | --- --- & \; x x --- & x x --- & --- \\
581 & --- x x & --- --- & \; --- --- & \; x x --- & | --- --- & \; --- --- & x x --- & --- \\
582 & --- --- & x x --- & \; x x --- & | x x --- & \; --- --- & \; --- --- & x x --- & x \\
583 & x x --- & --- --- & | x x --- & \; --- --- & |  --- --- & | x x --- & x x --- & --- \\
584 & --- x x & --- --- & | --- x x & \; --- --- & --- --- & | x x --- & x x --- & --- \\
585 & --- --- & x x --- & | --- x x & \; --- --- & --- --- & | --- ---\; & x x --- & x \\
586 & x x --- & x x --- & | --- ---\; & \;--- --- & --- --- & | --- ---\; & x x --- & --- \\
\hline
594 & --- --- & --- --- & \;--- --- & \; --- --- & \cellcolor{darkergray} | x x --- & | --- ---\; & x x --- & \multicolumn{1}{c}{x} \\
\end{tabular}

\vspace{14pt}
\textbf{Scansion of Penelope's request to the suitors (\textit{Odyssey} ii.96--102)}

\vspace{6pt}
\begin{tabular}{r | c c c c >{\columncolor{graycol}}c c}
96 & --- x x & --- --- & --- x $\mid$ x & --- x x & --- x x & --- x x \\
97 & --- x x & --- x x & --- $\mid$ x x & --- x x & $\mid$ --- x x & --- x \\
98 & --- x x & --- $\mid$ --- & --- $\mid$ x x & --- x x & --- x x & --- x \\
99 & --- --- & --- --- & --- x $\mid$ x & --- x x & $\mid$ --- x x & --- x \\
100 & --- x x & --- x x & --- x $\mid$ x & --- x x & --- x x & --- x \\
101 & --- --- & --- x x & --- x $\mid$ x & --- x x & --- x x & --- --- \\
102 & --- x x & --- --- & --- $\mid$ --- & --- --- & --- x x & --- --- \\
\end{tabular}
\end{singlespacing}
\end{table}


Even if coincidental, the similarities of the two passages, sans meter, are already marked---the two heroines operate in analogous settings (one of which is marked as that setting) and who have weaving-themed speeches that are the defining programs of both heroine's rebellious efforts. The two speeches also share an internally marked set of metrical consistencies (i.e., ones not demanded by the meter and not elsewhere internally followed so strictly) in (two variations of) the same location, which might make coincidence somewhat less likely.

If the argument from meter is valid, it would also seem significant that Lysistrata's speech specifically doesn't display the dactylic syllable structure as the link to Penelope's speech, even though that type of foot is available to that meter.\footnote{I don't want to be accused of special pleading here, so, to be clear, my claim here is not that the use of non-dactyl forms in \textit{Lysistrata} would be further evidence of a textual link with the \textit{Odyssey}. Rather, the claim is that the link is established before making note of non-dactyl feet, and that the non-dactyl feet then provide an interesting twist on the link.} Rather, the first ``linking'' set in Lysistrata's speech (every fifth foot) uses the syllable structure available to both epic meter and anapestic tetrameter, the spondee. The second ``linking'' set (every seventh foot) uses the reverse of the epic meter (and which is not available to epic meter).

\section{A Structural-Thematic Link}

Regardless of the felicity of the metrical argument, the two plays still seem to be connected by a structural-thematic link. The crux of Penelope's famous guile is her act of weaving (ὑφαίνεσκεν, ii.104) and subsequent loosening of (ἀλλύεσκεν, ii.105) her \textit{pharos}. The implication of the text (especially given the iterative forms of the verbs) is that Penelope's weaving-loosening trick would go on indefinitely if not for the intervention of one of the maids.\footnote{``[S]he transforms her work into endless weaving, a trait which the Greek imagination equates with a refusal or the impossibility of marriage, into the excessive prolonging of virginity'', Papadopoulou-Belmehdi, 44.} So, to reiterate, the backbone of Penelope's ruse, and thus to a large degree of her import and influence, is a constant (and potentially unending) ὑφαίνω-ἀναλύω, weaving-loosening, a back-and-forth structure.\footnote{Plato also highlights this particular unending back-and-forth aspect of Penelope's story. At 84a in the \textit{Phaedo}, Socrates says that the philosopher's soul wouldn't think it necessary ``to perform the endless task of Penelope working backwards some loom (ἀνήνυτον ἔργον πράττειν Πηνελόπης τινὰ ἐναντίως ἱστὸν μεταχειριζομένης). Plato, \textit{Phaedo}, trans. Chris Emlyn-Jones and William Preddy, LCL 36 (Cambridge, MA: Harvard University Press, 2017).}

When talking with the Proboulos, Lysistrata also frames her designs for the city with acts of weaving and loosening. The relevant pericope here is (what I will call) the ``women's work'' unit, which starts with the Proboulos questioning Lysistrata (565) and continues until Lysistrata takes up the Proboulos' subject change (590).\footnote{I don't think this view weakens my already weak metrical claim, that the primary wool-working speech has its own metrical pattern, since the domains of analysis are separate (thematic versus metrical).} This passage is framed by nearly the same set of words as Penelope's actions: weaving (ὑφαίνω, 586) and loosening (διαλύω, 566, 569).\footnote{Διαλύω has more overt political connotations than ἀναλύω, which makes sense given the political context.} In this case, however, the order is reversed. Lysistrata's program begins with loosening and ends with weaving.\footnote{I will emphasize here, though, that Lysistrata's speech is mostly not about weaving per se, but about pre-weaving actions. This fact seems to have been often missed. E.g., Martin, ``Lemnian Women'', 86; Vaio, ``Theme'', 373-376;  Hawkins, ``Misanthrope'', 157-158.} This sequence from loosening to weaving is more precise than mere thematic similarity, even though at first glance the two sections are referring to different domains of the polis' affairs, namely, foreign (the war) versus domestic (corruption/unity) affairs.\footnote{Though, in this case, the former might reasonably be considered a subset of the latter. Henderson, ``Commentary'', 141 blurs the two together too strongly.}

The sequence starts with the Proboulos' question (565-566) to Lysistrata as to how she plans to ``loosen'' (διαλύω) the problems pertaining to the war. Lysistrata then (567-570) takes up this question, claiming that she will ``loosen'' (διαλύω, 569) the war by treating it like messy ``spindled thread'' (κλωστήρ, 567)\footnote{Henderson's Loeb translation has ``ball of yarn'' for κλωστήρ. Aristophanes, \textit{Birds. Lysistrata. Women at the Thesmophoria}, trans. Jeffrey Henderson, LCL 179 (Cambridge, MA: Harvard University Press, 2000). I'm taking the \textit{LSJ} more literally here, \textit{LSJ, s.v.}, ``κλωστήρ'', last accessed November 2018, \url{http://www.perseus.tufts.edu/hopper/}.} and ``add in'' their ``spindles'' (ἄτρακτος, 568) to disentangle the thread. The Proboulos then (571-572) mocks the ``spindled thread'' (κλωστήρ) and ``spindles'' (ἄτρακτος) and also mocks (by addition) ``wool'' (ἔριον). Even if the idea of wool, as thread's material origin, has been implicit, it is still a different object than thread, one that is precisely ``backwards'' in the development of that thread. Lysistrata's response takes up only the ``wool'' (ἔριον, 573) added by the Proboulos. Each term is said twice and each uniquely said by each interlocutor, and, moreover, the original thread-spindle-wool sequence is repeated in the same order. The sequence can be diagrammed as follows:


\begin{table}[htbp]
\centering
\begin{singlespacing}
\begin{tabular}{c| c | c | c| c | c}
\multicolumn{2}{c|}{Lysistrata} & \multicolumn{3}{c|}{Proboulos} & Lysistrata \\
\hline
spindled thread & spindle & wool & spindled thread & spindle & wool \\
\end{tabular}
\end{singlespacing}
\end{table}


Thus, the transition from foreign to domestic policy is controlled by a well-ordered sequence within the larger wool theme. Rather than the shift from disentangling thread to cleaning wool being two disjoint examples of women's work, and metaphorical for mutually distinct domains, the thread is ``weaved'' through the two interlocutors until it becomes wool again.\footnote{Nor should the fact be lost, with respect to feminist readings of the text, that Lysistrata's plans are in each instance provoked by and implicit in words first said by the Proboulos, which seems like a way the text subtly strips agency/originality from Lysistrata. This fact could easily be combined with the other fact that Lysistrata, even before here speech, gives away her sewing materials to the Proboulos, further stripping her agency. Gwendolyn Compton-Engle, \textit{Costume in the Comedies of Aristophanes} (Cambridge: Cambridge University Press, 2015), 52 oddly claims that ``the well-known scene in which the women interact with the Proboulos (an Athenian magistrate) shows them physically in control of both garments and bodies'', even though the women seem to revoke this power.} Just like the supposedly dead Proboulos later fails to get his projected shroud, the ``loosening'' Lysistrata describes fails, on a textual-sequence level, to achieve the desired product, straightened thread, and instead ends up with the bare material substrate.

Continuing the sequence, Lysistrata's main wool-working speech starts with ``raw wool'' (πόκος), the furthest possible move ``backward'' for the more ambiguous ``wool'' (ἔριον). From there, Lysistrata in detail describes cleaning and preparing the wool for domestic use, and this description initially follows a natural, logical sequence (574-584): shearing (implicit from πόκος), washing (ἐκπλύνω), beating (ἐκραβδίζω), plucking (ἀπολέγω), unknotting (implicitly in the badness of πιλέω), carding through (διαξαίνω), carding (ξαίνω), mixing in everything ((εγ)καταμίγνυμι), and leading and collecting everything together (ξυνάγω καὶ συναθροίζω).\footnote{A slight disjoint already occurs here. The object of leading/collecting everything (ξυνάγω καὶ συναθροίζω) is ``flocks of (worsted) wool'' (κάταγμα), i.e., combed wool. Unlike all the previous stages, this one appears in the speech as a noun, without corresponding verb, i.e., without description of the wool actually going through the combing process. Moreover, the first instance of κάταγμα refers to Athens' colonies, which, as wool, are separate from the primary ``Athenian proper'' wool; i.e., the colony flocks implicitly haven't gone through the foregoing cleaning-carding sequence.}

At this point a problem occurs in the sequence. Lysistrata says that, after combining all of these ``flocks'' of (worsted) wool together, they should be ``made into a big clew of wool'' (ποιήσαι τολύπην μεγάλην, 585-586). This action, however, is, strictly speaking, impossible, because this wool hasn't been spun into fibers yet,\footnote{``From Sheep to Rug'', Nejad Rugs, modified 2007, \url{https://www.nejad.com/consumer/sheep\_to\_rug.htm}.} and, in a speech in which the processes have been sequentially and explicitly described, it would be infelicitous to claim that the (relatively important) spinning process is being shuffled in with the yarn-combining process.\footnote{The impossibility of ``making'' raw wool from non-raw wool, might be invoked here. I think there's enough ambiguity in πόκος to provide slippage between raw and non-raw wool. In any case, though, I'm privileging the explicit sequence over the literal material possibility, a privileging which seems condoned by the text's own tight sequencing of steps. What's important is that a step is skipped between the flocks and the yarn. This is not the case between wool and raw wool, even if it's not literally possible.} The sequence is broken right at the point where a usable product is made. As with foregoing failures of production, this break further calls into question the efficacy of Lysistrata's designs, and therefore of women's work more generally, to produce the \textit{chlaina}, along with the garment's democratic metaphorical implications.

The skipped spinning step is one in which wool is made into thread, a process which requires spindles. The ``women's work'' unit, as already seen, begins precisely with these two materials (κλωστήρ, 567, 571; ἄκτρακτος, 568, 571). The natural implication might be to ``fill in'' the missing step with the thread sequence previously described. This implication can be further supported in two ways. First, the gathering/collecting is preceded by the act of ``taking'' all the mixed up flocks of wool (τὸ κάταγμα λαβόντας, 585), and this semantically weak verb could refer to its use at the beginning of the women's work unit, when Lysistrata makes the first step ``taking [the spindled thread] in this way'' (ὧδε [αὐτὸν] λαβοῦσαι, 568). The weak verb acts as a link between the first instance and the second, and the demonstrative, then, refers not only to the stage actions,\footnote{Henderson, ``Commentary'', 142.} but proleptically to the gap in the sequence. More suggestively, the verbs directly preceding the this gap, ``leading and gathering together'' (ξυνάγω καὶ συναθροίζω), have strong military connotations, as if the wool being gathered are soldiers. The war, soldiers, is exactly what the spindles were invoked to fix at the beginning of the women's work unit.

To recap, the women's work sequence starts with thread and spindles, regresses to raw wool, and then progresses until the gap at the end. By filling in the sequence's later gap with the prior thread/spindle information, the implication is that the recently gathered together (συναθροίζω) flocks (κάταγμα) will again be turned back into raw wool, rather than be made into beautiful and inclusive yarn and \textit{chlaina}. The sequence, then, becomes a recursive back-and-forth sequence between raw wool and combed wool, and the actual wool products are separate and false conclusions. The yarn and \textit{chlaina} remain disconnected from the work and therefore unproduced.\footnote{Which is marked compared to women's ``traditional role as providers of textiles'', Compton-Engle, \textit{Costume}, 12.} Just like Penelope's weaving and loosening, Lysistrata's polis program presents a potentially unending back-and forth structure.

\section{Scene-Internal Political Implications}

In several ways, then, if subtly, the wool-working discussion highlights the necessary ineffectualness of Lysistrata's plan. Her designs fail to allow the Proboulos to make his expected shroud, and thread becomes raw wool instead of straightened thread. The yarn and \textit{chlaina} have no means of being produced, since, structurally speaking, once the wool gets to the yarn-spinning stage, the regression back to raw wool restarts. The back-and-forth sequence in \textit{Lysistrata} can be generalized in two equivalent ways:


\begin{table}[htbp]
\centering
\begin{singlespacing}
\textbf{Wool sequences in \textit{Lysistrata} (565--586)}

\vspace{6pt}
\footnotesize
\setlength{\dashlinedash}{1.5pt}
\setlength{\dashlinegap}{1.5pt}
\begin{tabular}{c:c:c:c:c:c:c}
%\hdashline
& \textbf{1} & \textbf{2} & \textbf{3} & \textbf{4} & \textbf{5} & \textbf{6 = 1} \\
\hdashline
\textbf{textual} & bunched & \multirow{2}{*}{wool} & \multirow{2}{*}{raw wool} & cleaned/ & bunched & bunched \\
\textbf{sequence} & thread & & & combed wool & combed wool & thread \\
\hdashline
\textbf{generalized} & ``clumped'' & ``unclumped'' & ``clumped'' & ``unclumped'' & ``clumped'' & ``clumped'' \\
\textbf{sequence 1} & wool & wool & wool & wool & wool & wool \\
\hline
\textbf{natural} & \multirow{2}{*}{raw wool} & cleaned/ & bunched & \multirow{2}{*}{bunched thread} & \multirow{2}{*}{wool} & \multirow{2}{*}{raw wool} \\
\textbf{sequence} & & combed wool & combed wool & & & \\
\hdashline
\textbf{generalized} & ``clumped'' & ``unclumped'' & ``clumped'' & ``clumped'' & ``unclumped'' & ``clumped'' \\
\textbf{sequence 2} & wool & wool & wool & wool & wool & wool \\
%\hdashline
\end{tabular}
\end{singlespacing}
\end{table}


In other words, the back-and-forth sequence can be generalized as a sequence between various modes of ``clumpy'' messes of wool and various modes of ``non-clumpy'' orders of wool. This characterization doesn't add information to the raw-combed wool sequence, but it highlights the political implications of the failures that occur in this part of the text. Namely, as the text makes explicit, the various stages of wool refer to different inhabitants of the polis, everyone from state officials to non-citizens. The supposedly ``democratic'' acts of this sequence are those of combining different kinds of polis-dwellers into one ``basket of unity'' (καλαθίσκον κοινὴν εὔνοιαν, 579\footnote{Perhaps a periphrasis for the \textit{Odyssey}'s \textit{homonoia}. Cf. Vetter, \textit{``Women's Work''}, 33-61.}).\footnote{Major, ``Court'', 138 claims Lysistrata ``call[s] for listening to and embracing the broadest possible coalition'' and desires that coalition to participate in the ``deliberative process''.} The democratic reforms, though, by being linked to the unending back-and-forth of the wool sequence, is itself implied to be stuck in a clump-unclump back-and-forth, cordoned off from the final, beautiful woven product. Put simply, democratic reforms are presented here as inefficacious because heterogenous ``clumps'' of people can't produce anything useful or aesthetic.\footnote{Contra Henderson's claim that ``this [wool-working] metaphor (anticipated at 519-20, 532-8) \dots  illustrates how the wives can offer, in terms of their own sphere, useful advice to the city'', Henderson, ``Commentary'', 141. The anticipatory metaphors he cites also support the inefficacy of women's work.}

The inefficacy of democratic reforms implied by the text isn't strictly dependent on its link to the Penelope story. The link does, however, provide an ironic relief for the failure of Lysistrata's plan. The aristocratic Penelope who is, among many superlatives, ``clever above all women'' (ἥ τοι πέρι κέρδεα οἶδεν, ii.88) and who is, as has been shown, skilled enough to actually produce something of value. The ostensibly (pro-)democratic Lysistrata, however, is not described in such superlative terms and cannot produce the desired product.\footnote{Women in the play are, e.g., ``knavish'' (πανοῦργος, 12).} Thus, maintaining the Penelope-Lysistrata link provides an aristocratic counterpoint to the democratic failures.\footnote{Even the syntax of the two speeches reflect this. Penelope's speech is a tightly controlled set of subordinate clauses, while Lysistrata's speech is mostly a ``clumped'' series of conjunctions and infinitives.} Having the counterpoint better lends the scene in \textit{Lysistrata} to a comedic interpretation, because it provides a canonical example by which the women in \textit{Lysistrata} might have ``known better'' and against which they may be denigrated.

\section{Political Implications of \textit{Lysistrata}, Briefly}

In the similar but slightly larger domain of clothes more generally, the inefficacy of democracy seems paralleled in \textit{Lysistrata}. The obvious parallel to the back-and-forth structure of the wool-working section is the dis- and rerobing of both semi-choruses. This back-and-forth, isn't strictly marked, as choruses routinely strip in Old Comedy.\footnote{Compton-Engle, \textit{Costume}, 53.} However, ``the orchestra [being] now filled with stage-naked old men and old women was a comic spectacle rivaling in its own way that of Old Comedy's most elaborately costumed choruses \dots  [and was] a bold and apparently unprecedented move''.\footnote{Ibid., 54. Moreover, the Women's Chorus Leader's command ``to shake off their old age (ἀποσείσασθαι τὸ γῆρας τόδε, 670)'' is an ``impossible'' act analogous to the ``impossible'' move from non-raw to raw wool.} Once again, the condition in which people don't have the standard covering/clothing, clothing ostensibly made by women, is highlighted. Similarly, while Lysistrata claims her seduction plan requires fancy clothes (42-48), the actual seduction scene with Myrrhine and Kinesias (864-953) instead focuses on tactile tantalization, as if fancy clothes (even if the actor wore them) wouldn't have had enough seductive power. In both cases, the potential use value of women's products are shirked in favor of women's tactile, naked, qualities.

More importantly, the resolution of the war is implicated in the failure of clothing/weaving. Reconciliation is literally a naked woman, and, in her function as a divided up map (at least a superset of Athens and the demos), her uncovered body recalls Lysistrata's failure, her inability, to produce a demos-covering \textit{chlaina} for her. In fact a second, and the only other, \textit{chlaina}\footnote{Also mentioned by Compton-Engle, \textit{Costumes}, 56.} is mentioned by Lysistrata as covering cold Athenians. This \textit{chlaina}, though, exists only in the memory of a past time and is, crucially, a product of the oligarchic Spartans.\footnote{The same oligarchs who get the last words in the play.}

The negotiations resemble nothing like the supposedly democratic plan envisioned by Lysistrata. The heroine leads the accords without input from the metics or other ``lower'' citizens, not even the other women.\footnote{Relatedly, at one point, Lysistrata's disdain for maidens (κόρη, 593) sitting home alone, directly contradicts the Women's Chorus Leaders' preference for sitting at home like a maiden (κόρη, 473). It's also interesting to note that both Penelope's and Lysistrata's unending sequences are ended by outside forces. In Penelope's case, Penelope can maintain her aloofness and cleverness, since she is betrayed. Lysistrata, though, in taking somewhat despotic control, acts as her own betrayer, further ironizing the purported democratic inclinations.} Nor do the negotiations resemble some other clear political organization, but, instead, the negotiations themselves seem generally farcical. The negotiations ``focus on this happy end-product [sex] \dots  rather than the issues'',\footnote{Henderson, ``Commentary'', 196.} and ``result in no concrete political treaties''.\footnote{Vetter, \textit{``Women's Work}, 74.} The negotiations themselves, and therefore their falseness,\footnote{As well as the obvious comedic intent in depicting sex-crazed men with erect penises.} are predicated on the nakedness of Reconciliation. By extension, the falseness of the negotiations are predicated on the failure of Lysistrata's plan to produce a garment and democratic unity. This doesn't however, imply that a successful instantiation of democratic procedures would save the negotiations, let alone the polis. Democracy has already been shown to be stuck in an ineffective clumping-unclumping loop, linked to the ineffective wool-working loop, and nothing has swooped in to save it. Rather, democracy is an irrelevant structure for negotiation, and nothing (the) women could do could change that.

\section{Conclusion}

Ultimately, \textit{Lysistrata} seems to provide a condemnatory perspective on both democratic forms of governance and women's potential emancipatory roles in the demos. Both democracy and women's work are structures that get caught in endless loops of inefficacy. The implications here for the genre of comedy might also be significant. On the one hand, comedy, here, is portrayed as somewhat parasitic upon the Penelope tradition for its potential positive alternative (the aristocracy) to the failed democratic reforms. On the other hand, Lysistrata's wool-working speech is placed in the center of the play and acts as the link (\dots  the ``thread'') between the two distinct plots of the play. If weaving is, for the play, a systematically unaccomplished, impossible activity, the implication would be that the comedy itself, and, by extension, perhaps the genre of comedy at large, might fail to accomplish anything. Instead, comedy becomes another raucous meeting of sex-craved bureaucrats.

\section{References}
%\addcontentsline{toc}{section}{References}
\begin{singlespacing}
\begin{hangparas}{.25in}{1}
Aristophanes. \textit{Birds. Lysistrata. Women at the Thesmophoria}. Translated by Jeffrey Henderson. LCL 179. Cambridge, MA: Harvard University Press, 2000.

---. \textit{Lysistrata}. Edited by Jeffrey Henderson. Oxford: Clarendon Press, 1987.

Compton-Engle, Gwendolyn. \textit{Costume in the Comedies of Aristophanes}. Cambridge: Cambridge University Press, 2015.

Dillon, Matthew. ``The \textit{Lysistrata} as a Post-Deceleian Peace Play''. \textit{Transactions of the American Philological Association} 117 (1987): 97--104.

Faraone, Christopher A. ``Salvation and Female Heroics in the Parodos of Aristophanes' \textit{Lysistrata}''. \textit{The Journal of Hellenic Studies} 117 (1997): 38--59.

Foley, Helene P. ``The `Female Intruder' Reconsidered: Women in Aristophanes' \textit{Lysistrata} and \textit{Ecclesiazusae}''. \textit{Classical Philology} 77, no.~1 (Jan. 1982): 1--21.

Hawkins, Tom. ``Seducing a Misanthrope: Timon the Philogynist in Aristophanes' \textit{Lysistrata}''. \textit{Greek, Roman, and Byzantine Studies} 42 (2001): 143--162.

Henderson, Jeffrey. ``Commentary''. In \textit{Lysistrata}, edited by Jeffrey Henderson, 65--222. Oxford: Clarendon Press, 1987.

---. ``Introduction''. In \textit{Lysistrata}, edited by Jeffrey Henderson, xv--lxxii. Oxford: Clarendon Press, 1987.

Homer. \textit{Odyssey}. Volumes I--II. Translated by A.T. Murray. LCL 104--105. Cambridge, MA: Harvard University Press, 1919.

\textit{LSJ}. Last accessed November 2018. \url{http://www.perseus.tufts.edu/hopper/}.

Major, Wilfred E. \textit{The Court of Comedy: Aristophanes, Rhetoric and Democracy in Fifth-Century Athens}. Columbus: Ohio State University Press, 2013.

Martin, Richard P. ``Fire on the Mountain: `Lysistrata' and the Lemnian Women''. \textit{Classical Antiquity} 6, no.~1 (Apr. 1987): 77--105.

McClure, Laura. ``Courtesans Reconsidered: Women in Aristophanes' \textit{Lysistrata}''. \textit{Eugesta} 5 (2015): 54--84.

McGlew, James. \textit{Citizens on Stage: Comedy and Political Culture in the Athenian Democracy}. Ann Arbor: University of Michigan Press, 2002.

Nejad Rugs. ``From Sheep to Rug''. Copyright 2007. Last accessed November 2018. \url{https://www.nejad.com/consumer/sheep_to_rug.htm}.

Papadopoulou-Belmehdi, Ioanna. ``Greek Weaving or the Feminine in Antithesis''. \textit{Diogenes} 42/3, no.~167 (Fall 1994): 39--56.

Plato. \textit{Phaedo}. Translated by Chris Emlyn-Jones and William Preddy. LCL 36. Cambridge, MA: Harvard University Press, 2017.

Stroup, Sarah Culpepper. ``Designing Women: Aristophanes' \textit{Lysistrata} and the `Hetairization' of the Greek Wife''. \textit{Arethusa} 37, no.~1 (Winter 2004): 37--73.

Vaio, John. ``The Manipulation of Theme and Action in Aristophanes' \textit{Lysistrata}''. \textit{Greek, Roman, and Byzantine Studies} 14 (Winter 1973): 369--380.

Vetter, Lisa Pace. \textit{``Women's Work'' as Political Art}. Lanham, MD: Lexington Books, 2005.

Westlake, H.D. ``The `Lysistrata' and the War''. \textit{Phoenix} 34, no.~1 (Spring 1980): 38--54.
\end{hangparas}
\end{singlespacing}

\end{document}
